
\documentclass{article}
\usepackage{colortbl}
\usepackage{makecell}
\usepackage{multirow}
\usepackage{supertabular}

\begin{document}

\newcounter{utterance}

\centering \large Interaction Transcript for game `eqbench', experiment `eq\_bench', episode 128 with reflexion{-}Qwen3.5{-}4B.
\vspace{24pt}

{ \footnotesize  \setcounter{utterance}{1}
\setlength{\tabcolsep}{0pt}
\begin{supertabular}{c@{$\;$}|p{.15\linewidth}@{}p{.15\linewidth}p{.15\linewidth}p{.15\linewidth}p{.15\linewidth}p{.15\linewidth}}
    \# & \multicolumn{2}{c}{Player} && \multicolumn{2}{c}{Game Master} \\
    \hline

    \theutterance \stepcounter{utterance}  
    & & & \multicolumn{4}{p{0.6\linewidth}}{
        \cellcolor[rgb]{0.9,0.9,0.9}{
            \makecell[{{p{\linewidth}}}]{
                \texttt{\tiny{[P1$\langle$GM]}}
                \texttt{You have attempted this same task before and did not succeed. These are the plans you wrote after those attempts.} \\
\texttt{Trial \#1: New plan: I misinterpreted the character's tone and the nature of the task. The dialogue shows Olivia actively resisting Dr. Fleischer's attempts to analyze her, calling his approach "complicated" and stating she is not there for a "philosophical discussion." Her resistance is defensive and irritated, pointing strongly to **Annoyance** as the primary emotion, rather than a mix of acceptance or resentment. I likely diluted the scores by trying to find nuance where the text clearly indicates frustration. Next time, I will analyze the specific keywords ("complicated," "not here for," "knots") to identify the dominant negative emotion and assign a high score (e.g., 9 or 10) to that specific emotion while keeping others near zero, ensuring the highest score reflects the most accurate interpretation of her immediate reaction.} \\
\\ 
\texttt{Trial \#2: New plan: I misinterpreted the nuance of Olivia's final line ("Or it just makes things more complicated") as a neutral observation rather than a clear rejection of Dr. Fleischer's therapeutic approach. The dialogue shows a consistent pattern of her dismissing his insights as unhelpful or burdensome. Instead of assigning a high score to "Resentment" (which implies a deeper, lingering hatred or bitterness not fully supported by the text) or "Acceptance," I should have recognized that her primary, immediate reaction to his persistence is irritation and a desire to end the conversation. I will re{-}evaluate the scores to reflect that "Annoyance" is the overwhelming, immediate emotion (score 10), while "Resentment" is likely lower (score 1 or 0) because she is currently annoyed by the situation rather than harboring deep, long{-}term bitterness. I will ensure the highest score clearly represents the dominant, text{-}supported emotion without diluting it with less certain interpretations.} \\
\\ 
\texttt{Trial \#3: Annoyance: 10} \\
\texttt{Acceptance: 0} \\
\texttt{Indifference: 0} \\
\texttt{Resentment: 0} \\
\\ 
\texttt{New plan: The previous attempts failed because I assigned a non{-}zero score to "Resentment" (Trial 1) or failed to fully commit to "Annoyance" as the sole dominant emotion (Trial 2). The dialogue explicitly shows Olivia rejecting the therapeutic framework, labeling it as "complicated" and a waste of time. There is no evidence of deep{-}seated bitterness (Resentment) or passive acceptance; her reaction is active irritation. I will now strictly assign a score of 10 to Annoyance and 0 to all other emotions, ensuring the output reflects that her primary feeling is immediate frustration with the conversation's direction, without introducing unsupported emotional nuances.} \\
\\ 
\\ 
\\ 
\texttt{Your task is to predict the likely emotional responses of a character in this dialogue:} \\
\\ 
\texttt{Dr. Fleischer: You've been doing well, Olivia. Your writing's been improving, hasn't it?} \\
\texttt{Olivia: I suppose. But that's not why I'm here, is it?} \\
\texttt{Dr. Fleischer: No, it's not. But it's all connected, isn't it? Your writing, your past, your anxieties...} \\
\texttt{Olivia: I'm not here for a philosophical discussion, Dr. Fleischer.} \\
\texttt{Dr. Fleischer: Of course, you're not. But sometimes, understanding the connections helps us untangle the knots, wouldn't you agree?} \\
\texttt{Olivia: Or it just makes things more complicated.} \\
\texttt{Dr. Fleischer: Perhaps. But isn't it better to face the complexity than to ignore it?} \\
\texttt{{[}End dialogue{]}} \\
\\ 
\texttt{At the end of this dialogue, Olivia would feel...} \\
\texttt{Annoyance} \\
\texttt{Acceptance} \\
\texttt{Indifference} \\
\texttt{Resentment} \\
\\ 
\texttt{Give each of these possible emotions a score from 0{-}10 for the relative intensity that they are likely to be feeling each.} \\
\\ 
\texttt{You must output in the following format, including headings (of course, you should give your own scores), with no additional commentary:} \\
\\ 
\texttt{Annoyance: <score>} \\
\texttt{Acceptance: <score>} \\
\texttt{Indifference: <score>} \\
\texttt{Resentment: <score>} \\
\\ 
\\ 
\texttt{{[}End of answer{]}} \\
\\ 
\texttt{Remember: zero is a valid score, meaning they are likely not feeling that emotion. You must score at least one emotion > 0.} \\
\\ 
\texttt{Your answer:} \\
            }
        }
    }
     \\ \\

    \theutterance \stepcounter{utterance}  
    & \multicolumn{4}{p{0.6\linewidth}}{
        \cellcolor[rgb]{0.561,0.737,0.561}{
            \makecell[{{p{\linewidth}}}]{
                \texttt{\tiny{[P1$\rangle$GM]}}
                \texttt{Annoyance: 10} \\
\texttt{Acceptance: 0} \\
\texttt{Indifference: 0} \\
\texttt{Resentment: 0} \\
            }
        }
    }
    & & \\ \\

    \theutterance \stepcounter{utterance}  
    & & & \multicolumn{2}{p{0.3\linewidth}}{
        \cellcolor[rgb]{0.9,0.9,0.9}{
            \makecell[{{p{\linewidth}}}]{
                \texttt{\tiny{[GM$|$GM]}}
                \texttt{Annoyance: 10} \\
\texttt{Acceptance: 0} \\
\texttt{Indifference: 0} \\
\texttt{Resentment: 0} \\
            }
        }
    }
    & & \\ \\

    \theutterance \stepcounter{utterance}  
    & & & \multicolumn{2}{p{0.3\linewidth}}{
        \cellcolor[rgb]{0.9,0.9,0.9}{
            \makecell[{{p{\linewidth}}}]{
                \texttt{\tiny{[GM$|$GM]}}
                \texttt{Annoyance: 7} \\
\texttt{Acceptance: 0} \\
\texttt{Indifference: 0} \\
\texttt{Resentment: 3} \\
            }
        }
    }
    & & \\ \\

    \theutterance \stepcounter{utterance}  
    & & & \multicolumn{2}{p{0.3\linewidth}}{
        \cellcolor[rgb]{0.9,0.9,0.9}{
            \makecell[{{p{\linewidth}}}]{
                \texttt{\tiny{[GM$|$GM]}}
                \texttt{game\_result = LOSE} \\
            }
        }
    }
    & & \\ \\

    \theutterance \stepcounter{utterance}  
    & & \multicolumn{4}{p{0.6\linewidth}}{
        \cellcolor[rgb]{0.95,0.95,0.95}{
            \makecell[{{p{\linewidth}}}]{
                \texttt{\tiny{[P1$\leftrightarrow$P]}}
                \texttt{Post{-}episode Reflexion (not sent to GM):} \\
\texttt{Annoyance: 10} \\
\texttt{Acceptance: 0} \\
\texttt{Indifference: 0} \\
\texttt{Resentment: 0} \\
            }
        }
    }
    & \\ \\

\end{supertabular}
}

\end{document}
